\chapter{Hiện thực}

\section{Thiết lập môi trường phát triển}
\label{sec:dev_setup}

\subsection{Nền tảng Phần cứng}
Hệ thống được triển khai và kiểm thử trên bộ kít phát triển \textbf{Xilinx Kria KV260 Vision AI Starter Kit}. Đây là nền tảng chuyên dụng cho các ứng dụng thị giác máy tính tại biên (Edge AI).
\begin{itemize}
    \item \textbf{System-on-Module (SOM):} Kria K26 SOM chứa chip Zynq UltraScale+ MPSoC (XCK26).
    \item \textbf{Tài nguyên:} 4GB RAM DDR4 (dùng chung cho PS và PL), 4 nhân ARM Cortex-A53 và khối logic FPGA với 256K Logic Cells.
    \item \textbf{Ngoại vi:} Camera MIPI CSI-2 (Raspberry Pi Camera V2), cổng Ethernet 1G, và cổng DisplayPort.
\end{itemize}


\subsection{Bộ công cụ thiết kế của Xilinx (Xilinx Design Tools)}
Quy trình thiết kế phần cứng và hệ điều hành được thực hiện trên phiên bản 
sử dụng \textbf{2023.2} để đảm bảo tính tương thích:

\begin{enumerate}
    \item \textbf{Vivado Design Suite:} Sử dụng để thiết kế kiến trúc phần cứng 
    (Block Design), tổng hợp mạch (Synthesis) và tạo bitstream cấu hình cho FPGA.
    \item \textbf{Vitis Unified Software Platform:} Dùng để phát triển Firmware 
    bare-metal và các ứng dụng tăng tốc phần cứng.
    \item \textbf{PetaLinux Tools:} Bộ công cụ xây dựng hệ điều hành 
    Linux nhúng tùy biến.
\end{enumerate}

\subsection{Môi trường biên dịch RISC-V (RISC-V Toolchain)}
Khác với các vi xử lý RISC-V tiêu chuẩn, nhân PULP sử dụng tập 
lệnh mở rộng để tăng tốc DSP và vòng lặp. Do đó, hệ thống yêu cầu bộ biên dịch 
chuyên biệt \texttt{pulp-riscv-gnu-toolchain}.

\begin{itemize}
    \item \textbf{Compiler:} \texttt{riscv32-unknown-elf-gcc}.
    \item \textbf{Cấu hình Build:} Hỗ trợ kiến trúc \texttt{rv32imcxpulpv2} 
    (Tập lệnh cơ bản + Nhân + Nén + Mở rộng PULP).
    \item \textbf{Thư viện:} Newlib (thư viện C chuẩn rút gọn cho hệ thống nhúng).
\end{itemize}

%============================================================
%
%============================================================
\section{Quy trình triển khai hệ thống}
\label{sec:workflow}

Để hiện thực hóa một hệ thống SoC phức tạp kết hợp giữa phần cứng tùy biến (FPGA) 
và phần mềm nhúng (Embedded Software), quá trình triển khai được chia thành 
5 giai đoạn tuần tự theo mô hình thiết kế dựa trên nền tảng (Platform-based Design). 

\subsection{Mô tả các giai đoạn}

Quá trình hiện thực được thực hiện lần lượt theo các bước sau:

\begin{enumerate}
    \item \textbf{Giai đoạn 1: Chuẩn bị Môi trường (Setup \& Toolchain)}
    \begin{itemize}
        \item Cài đặt và cấu hình các công cụ cốt lõi: Vivado (Phần cứng), Vitis (Phần mềm), PetaLinux (Hệ điều hành).
        \item Thiết lập chuỗi công cụ biên dịch chéo (Cross-compiler Toolchain) cho kiến trúc RISC-V PULP.
    \end{itemize}

    \item \textbf{Giai đoạn 2: Tối ưu hóa Mô hình AI (Model Optimization)}
    \begin{itemize}
        \item Huấn luyện và tinh chỉnh (Fine-tune) mạng MobileNetV1 trên máy tính cá nhân (Host PC).
        \item Thực hiện lượng tử hóa (Quantization) từ Float32 sang Int8.
        \item Xuất mô hình dưới dạng mảng dữ liệu C (Weights/Bias Header files) để sẵn sàng nhúng xuống phần cứng.
    \end{itemize}

    \item \textbf{Giai đoạn 3: Thiết kế Phần cứng SoC (Hardware Implementation)}
    \begin{itemize}
        \item Tích hợp IP Core RISC-V và khối tăng tốc CNN Accelerator vào nền tảng Zynq UltraScale+ trên Vivado.
        \item Thiết lập các kết nối AXI Interconnect, xung nhịp và bản đồ địa chỉ.
        \item Tổng hợp mạch (Synthesis), thực thi (Implementation) và tạo Bitstream (.bit).
    \end{itemize}

    \item \textbf{Giai đoạn 4: Phát triển Phần mềm Nhúng (Firmware \& OS)}
    \begin{itemize}
        \item Xây dựng hệ điều hành Linux tùy biến bằng PetaLinux, kích hoạt các driver cần thiết (UIO, CMA).
        \item Viết chương trình Firmware (Bare-metal) cho RISC-V để điều khiển khối Accelerator.
        \item Viết ứng dụng quản lý trên ARM (Python/C++) để giao tiếp và hiển thị kết quả.
    \end{itemize}

    \item \textbf{Giai đoạn 5: Tích hợp và Kiểm thử (Integration \& Verification)}
    \begin{itemize}
        \item Nạp Bitstream và khởi động hệ thống trên board Kria KV260.
        \item Chạy thực nghiệm với Video đầu vào từ Camera.
        \item Đo đạc thông số: Tốc độ (FPS), Độ chính xác (Accuracy) và Tài nguyên sử dụng.
    \end{itemize}
\end{enumerate}
%============================================================
%
%============================================================
\newpage
\section{Huấn luyện và tối ưu hóa Mô hình AI}

\subsection{Chuẩn bị Dữ liệu (Dataset)}
Hệ thống sử dụng bộ dữ liệu "Human Detection Dataset" từ Kaggle: 
\href{https://www.kaggle.com/datasets/jcoral02/inriaperson}{Kaggle - INRIA Person Dataset}
, bao gồm gần 300 ảnh 
ảnh có gán nhãn (Labeled) với hai lớp: \texttt{Human} và \texttt{Background}. 
Dữ liệu được chia theo tỷ lệ 80:20 cho huấn luyện và kiểm thử.

\subsection{Huấn luyện (Training Process)}
Quá trình huấn luyện sử dụng kỹ thuật Transfer Learning trên nền tảng TensorFlow/Keras:
\begin{itemize}
    \item \textbf{Base Model:} MobileNetV1 ($ \alpha=0.25 $) đã được pre-train trên 
    ImageNet. Các lớp tích chập được đóng băng (Freeze).
    \item \textbf{Top Layer:} Thêm lớp Fully Connected mới để phân loại 2 lớp.
    \item \textbf{Loss Function:} Binary Cross-Entropy.
\end{itemize}

% % Chèn đoạn code Python train mô hình (minh họa)
% \begin{lstlisting}[language=Python, caption=Cấu hình Transfer Learning cho MobileNet]
% base_model = tf.keras.applications.MobileNet(
%     input_shape=(224, 224, 3), alpha=0.25, include_top=False, weights='imagenet')
% base_model.trainable = False 

% model = tf.keras.Sequential([
%   base_model,
%   tf.keras.layers.GlobalAveragePooling2D(),
%   tf.keras.layers.Dense(1, activation='sigmoid')
% ])
% \end{lstlisting}

\subsection{Lượng tử hóa và Chuyển đổi (Quantization \& Conversion)}
Để triển khai mô hình trên vi xử lý RISC-V không có bộ xử lý số thực (FPU) mạnh mẽ, mô hình cần được chuyển đổi sang định dạng tính toán số nguyên 8-bit (INT8). Quy trình này sử dụng \texttt{TensorFlow Lite Converter}:

\begin{itemize}
    \item \textbf{Full Integer Quantization:} Chuyển đổi toàn bộ trọng số (Weights) và đầu ra các lớp (Activations) về khoảng giá trị [-128, 127].
    \item \textbf{Representative Dataset:} Sử dụng một tập con gồm 100 ảnh từ tập huấn luyện để bộ chuyển đổi hiệu chỉnh (calibrate) dải động (dynamic range) của các tham số.
\end{itemize}

% \begin{lstlisting}[language=Python, caption=Mã chuyển đổi TFLite và Lượng tử hóa]
% # Hàm generator cung cấp dữ liệu mẫu
% def representative_data_gen():
%     for input_value in tf_dataset.take(100):
%         yield [input_value]

% converter = tf.lite.TFLiteConverter.from_keras_model(model)
% converter.optimizations = [tf.lite.Optimize.DEFAULT]
% converter.representative_dataset = representative_data_gen
% # Ép kiểu tính toán và đầu ra về int8
% converter.target_spec.supported_ops = [tf.lite.OpsSet.TFLITE_BUILTINS_INT8]
% converter.inference_input_type = tf.int8
% converter.inference_output_type = tf.int8

% tflite_model_quant = converter.convert()
% # Lưu ra file .tflite
% with open('mobilenet_quant.tflite', 'wb') as f:
%     f.write(tflite_model_quant)
% \end{lstlisting}

\subsection{Xuất dữ liệu sang C Header (Export to C)}
Bước cuối cùng của Phase 1 là chuyển đổi file nhị phân \texttt{.tflite} 
thành các mảng hằng số ngôn ngữ C để nhúng vào Firmware của RISC-V. 
\begin{itemize}
    \item File \texttt{weights.h}: Chứa mảng trọng số.
    \item File \texttt{bias.h}: Chứa các hệ số chệch.
\end{itemize}

Dưới đây là cấu trúc dữ liệu sau khi chuyển đổi, sẵn sàng để nạp vào bộ nhớ FPGA:

% \begin{lstlisting}[language=C, caption=Trích đoạn file weights.h sau khi xuất]
% // File: weights.h
% // MobileNet Quantized Weights
% #include <stdint.h>

% const uint8_t conv1_weights[] __attribute__((aligned(16))) = {
%     0x04, 0xFE, 0x1A, 0x22, ... // Dữ liệu Hex dạng int8
% };

% const int32_t conv1_bias[] = {
%     120, -45, 300, ... // Bias thường lưu ở int32
% };
% \end{lstlisting}

\newpage
\subsection{Kết quả Huấn luyện ban đầu}
Kết quả thực nghiệm trên tập kiểm thử (Test set) cho thấy mô hình sau khi lượng tử hóa vẫn giữ được độ chính xác chấp nhận được cho bài toán phát hiện người:

\begin{table}[!h]
\centering
\begin{tabular}{|l|c|c|c|}
\hline
\textbf{Phiên bản Mô hình} & \textbf{Độ chính xác (Acc)} & \textbf{Kích thước} & \textbf{Ghi chú} \\ \hline
MobileNet (Float32) & 94.5\% & 16.2 MB & Chạy trên PC \\ \hline
MobileNet (Int8) & 92.8\% & 4.1 MB & Giảm 4x kích thước \\ \hline
\end{tabular}
\caption{So sánh hiệu năng mô hình trước và sau lượng tử hóa.}
\label{tab:model_result}
\end{table}
